\section{Full Pairing Proof for Emulated Fields}
\label{sec:emulated}

When the prover's native field $\F_n$ differs from BN254's $\Fq$ (e.g.\ verifying BN254 pairings inside a Plonk circuit over a different curve), all $\Fq$ arithmetic must be emulated via multi-limb representations in $\F_n$. This section first summarizes gnark's deferred multiplication approach~\cite{gnark}, then describes our contribution: batching those deferred checks into a single evaluation consisting only of flat $\F_n$ multiplications and additions.

\subsection{Deferred Emulated Multiplication in Gnark}

An element $a \in \F_e$ is represented by $\mathbf{n}$ limbs of $\mathbf{b}$ bits each in $\F_n$:
$$a = \sum_{i=0}^{\mathbf{n}-1} a_i \cdot 2^{i \cdot \mathbf{b}}, \quad a_i \in \F_n$$

The limb width satisfies $2\mathbf{b} < \log_2 p_n$ so that a product of two limbs fits in $\F_n$.

For a multiplication $a \cdot b \equiv r \pmod{p_e}$, the integer identity $a \cdot b = r + k \cdot p_e$ lifts to the \emph{carry-corrected polynomial check}~\cite{gnark}:

\begin{equation}
  \label{eq:emul-mul}
  a(X) \cdot b(X) = r(X) + k(X) \cdot p_e(X) + (2^{\mathbf{b}} - X) \cdot c(X)
\end{equation}

where $a(X) = \sum_i a_i X^i$ is the limb polynomial, and $c(X)$ is a \emph{carry polynomial}. The term $(2^{\mathbf{b}} - X) \cdot c(X)$ encodes the schoolbook carry chain: it subtracts $2^{\mathbf{b}} c_i$ from coefficient $i$ (removing overflow) and adds $c_i$ to coefficient $i+1$ (propagating it to the next limb). The elements $a$ and $b$ \emph{need not be reduced} modulo $p_e$ beforehand---they only need to be limb-bounded---which is what makes deferred verification possible.

Gnark's approach~\cite{gnark}: at each multiplication call site, compute $r, k, c$ via hint and return $r$ immediately; store the tuple $(a, b, r, k, c)$ for a later batch check. This defers all emulated multiplication costs to a single verification phase, avoiding redundant reduction steps.

\subsection{Our Contribution: RLC Batching into a Flat Check}

Gnark batches its deferred checks using a single-round challenge. Our contribution is to integrate this into the hierarchical IOP from Section~\ref{sec:iop}, and observe that the resulting batched check reduces entirely to \textbf{flat $\F_n$ multiplications and additions}---no limb arithmetic, no carries, no emulated operations whatsoever.

Given $m$ deferred tuples $\{(a_i, b_i, r_i, k_i, c_i)\}$ and a challenge $z_2 \xleftarrow{\$} \F_n$, we form accumulated polynomials:
$$K_{\text{acc}}(X) = \sum_{i=0}^{m-1} z_2^i \cdot k_i(X), \qquad C_{\text{acc}}(X) = \sum_{i=0}^{m-1} z_2^i \cdot c_i(X)$$

The verifier samples $x_2 \xleftarrow{\$} \F_n$ and checks:

\begin{equation}
  \label{eq:emul-batched}
  \underbrace{\sum_{i=0}^{m-1} z_2^i \bigl[a_i(x_2)\cdot b_i(x_2) - r_i(x_2)\bigr]}_{\text{LHS: flat } \F_n \text{ muls \& adds}} = \underbrace{K_{\text{acc}}(x_2) \cdot p_e(x_2) + (2^{\mathbf{b}} - x_2) \cdot C_{\text{acc}}(x_2)}_{\text{RHS: flat } \F_n \text{ muls \& adds}}
\end{equation}

This is the RLC of Eq.~\ref{eq:emul-mul} evaluated at $x_2$. Every term---the limb polynomial evaluations $a_i(x_2)$, $b_i(x_2)$, $r_i(x_2)$, $k_i(x_2)$, $c_i(x_2)$, the modulus evaluation $p_e(x_2)$, and the scalar $(2^{\mathbf{b}} - x_2)$---is a \textbf{flat scalar in $\F_n$}. The entire check is a sequence of $\F_n$ multiplications and additions; there is no emulated field arithmetic, no carry handling, and no limb-level bookkeeping at verification time. This is exactly the ``almost free'' check promised by the deferred approach.

\subsection{Hierarchical Four-Round IOP}

Combining with Section~\ref{sec:iop} gives a four-round protocol:

\textbf{Level 1 (Rounds I--II):} Pairing-level polynomial equations over $\F_e$, treated as exact (emulated correctness deferred). Prover commits remainders $\mathcal{R}_\text{pair}$, receives $z_1$, returns $Q_{\text{acc,pair}}$.

\textbf{Level 2 (Rounds III--IV):} After pairing-level evaluation point $x_1$ is fixed, prover commits the field-level remainder array $\mathcal{R}_\text{field}$ (the $r_i$ limb arrays), receives $z_2$, and returns $K_{\text{acc}}$ and $C_{\text{acc}}$. Verifier checks Eq.~\ref{eq:emul-batched} at a freshly sampled $x_2$.

\IOP{Four-Round Emulated Pairing IOP}{Verifier}{Prover}{
    \RDoes{Pairing computation with emulated $\F_e$; defers mul checks.\\ Prepares $\mathcal{R}_\text{pair}$, $\mathcal{Q}_\text{pair}$}{1}
    \RtoL{Round I: Commits pairing inputs and $\mathcal{R}_\text{pair}$}{2.5}
    \LtoR{Challenge $z_1 \xleftarrow{\$} \Fq$}{3.5}
    \RDoes{Computes $Q_{\text{acc,pair}}$ from $\mathcal{Q}_\text{pair}$ and $z_1$}{4.5}
    \RtoL{Round II: Commits $Q_{\text{acc,pair}}$}{5.5}
    \LDoes{Samples $x_1 \xleftarrow{\$} \Fq$; checks pairing-level identities}{6.5}
    \RtoL{Round III: Commits $\mathcal{R}_\text{field}$ (deferred mul remainders)}{8}
    \LtoR{Challenge $z_2 \xleftarrow{\$} \F_n$}{9}
    \RDoes{Computes $K_{\text{acc}}$, $C_{\text{acc}}$ from stored $k_i$, $c_i$ and $z_2$}{10}
    \RtoL{Round IV: Commits $K_{\text{acc}}$, $C_{\text{acc}}$}{11}
    \LDoes{Samples $x_2 \xleftarrow{\$} \F_n$; checks Eq.~\ref{eq:emul-batched} (flat $\F_n$)}{12}
    \LtoR{Accept/Reject}{13}
}

\subsection{Security and Fiat-Shamir}

\textbf{Soundness.} Each check (Eq.~\ref{eq:emul-mul}) is a polynomial identity of degree $2\mathbf{n}$ over $\F_n$. After RLC batching with challenge $(z_2, x_2)$, the bivariate Schwartz-Zippel lemma gives:
$$\delta_{s,\text{field}} \leq \frac{(m-1) + 2\mathbf{n}}{|\F_n|}$$
Combined with $\delta_{s,\text{pair}} \leq 2^{-245}$ (Section~\ref{sec:iop:soundness}), the total soundness $\delta_s \leq \delta_{s,\text{pair}} + \delta_{s,\text{field}}$ is negligible for $|\F_n| \approx 2^{254}$.

\textbf{Fiat-Shamir} makes this non-interactive with sequential hashes:
\begin{align*}
z_1 = H(\text{pairing inputs},\ \mathcal{R}_\text{pair}), \quad &
x_1 = H(z_1,\ Q_{\text{acc,pair}}) \\
z_2 = H(x_1,\ \mathcal{R}_\text{field}), \quad &
x_2 = H(z_2,\ K_{\text{acc}},\ C_{\text{acc}})
\end{align*}
Each challenge is bound to all prior transcript elements, preserving the temporal ordering that prevents retroactive adjustment of committed values.
